\chapter{Cast Care}

\section*{Definition and Overview}
Cast care is the set of practical steps taken to look after a limb that is immobilised in a plaster or fibreglass cast. Casts are used to hold broken bones in place while they heal, to protect a joint or repaired tendon after surgery, and to support severe sprains. A well-fitting cast keeps the injured part still so that healing can occur undisturbed, and most casts are worn for several weeks.

Proper cast care reduces the risk of complications such as pressure sores, skin irritation, and infection, and it keeps the cast strong and comfortable. It includes keeping the cast dry, managing swelling, protecting it from damage, exercising the joints that are not immobilised, and recognising the warning signs that mean the cast needs urgent review.

\section*{Epidemiology}
Broken bones are very common, and casting is one of the most frequently performed procedures in emergency departments and fracture clinics. Fractures of the wrist, forearm, and ankle are especially common in children, active younger adults, and older people who fall. Many people will wear a cast at some point in their lives, and for most the experience is straightforward when basic care rules are followed.

Complications of casting are also common enough to matter: skin irritation, pressure sores, and excessive swelling are the usual problems, while the rare but serious complication of compartment syndrome requires emergency treatment. Good cast care and early reporting of warning signs prevent most of these difficulties.

\section*{Aetiology and Pathophysiology}
The reason a cast is needed is usually an injury or operation that requires the injured part to be held still:

\begin{itemize}
    \item \textbf{Fractures:} bones must be kept aligned and motionless while the break heals
    \item \textbf{Dislocations and severe sprains:} damaged ligaments need rest and support
    \item \textbf{Tendon injuries:} repaired or strained tendons must be protected from pulling
    \item \textbf{After surgery:} to protect a repaired bone, joint, or soft tissue as it heals
    \item \textbf{Correction of deformity:} casts are occasionally used in children to gradually correct foot or limb position
\end{itemize}

A cast works by splinting the bones and joints and preventing movement that would disturb healing. Plaster casts are heavy and rigid; fibreglass casts are lighter and more durable. If the cast becomes wet, damaged, or too loose or tight, it no longer protects the injury properly and may itself cause harm by rubbing the skin, trapping moisture, or pressing on nerves and blood vessels.

\section*{Clinical Features}
The signs of a healthy, correctly fitting cast include a comfortable fit with some room for the fingers or toes to move, no severe pain, and normal colour, warmth, and feeling in the exposed hand or foot. Problems to watch for include:

\begin{itemize}
    \item Severe pain that is not relieved by pain medicine or by elevating the limb
    \item Swelling of the fingers or toes, or of the limb above or below the cast
    \item Numbness, tingling, or pins and needles in the hand or foot
    \item Fingers or toes that look pale, blue, or cold, or that cannot be moved
    \item A tight, throbbing, or "squeezing" feeling inside the cast
    \item Wet, soft, or cracked areas on the cast
    \item Itching, burning, or a strong smell coming from under the cast
    \item Redness or raw skin at the cast edges
\end{itemize}

\section*{Differential Diagnosis}
\begin{itemize}
    \item \textbf{Normal healing discomfort} vs.\ \textbf{compartment syndrome:} mild ache is expected, but severe pain out of proportion to the injury that is not relieved by elevation or analgesia is a warning sign
    \item \textbf{Expected swelling} vs.\ \textbf{a cast that is too tight} and cutting off circulation
    \item \textbf{Minor skin irritation} vs.\ \textbf{a pressure sore or infection} developing under the cast
    \item \textbf{Wet or softening cast} vs.\ \textbf{skin maceration or fungal infection} from trapped moisture
    \item \textbf{A loose cast after swelling settles} vs.\ \textbf{a broken or cracked cast} that no longer supports the bone
\end{itemize}

Any of these differences is best assessed by a health professional, who may need to inspect the skin by cutting a window in the cast or by replacing the cast altogether.

\section*{When to Seek Medical Care}
Contact the treating clinic or go to urgent care if the cast feels too tight, becomes soaked, cracks, or rubs painfully, or if there is a smell or discharge from beneath it. Fingers or toes that are swollen, numb, or blue need prompt assessment.

\textbf{Call 000 (or local emergency services) for:}
\begin{itemize}
    \item Severe, worsening pain in the injured limb that is not helped by elevation or pain relief
    \item Increasing tightness, burning, or a squeezing sensation inside the cast
    \item Fingers or toes that become pale, cold, blue, or completely numb
    \item Inability to move the fingers or toes
\end{itemize}

These features can indicate dangerously high pressure inside the limb and require urgent removal or splitting of the cast.

\section*{Diagnosis and Investigations}
\begin{itemize}
    \item \textbf{Clinical assessment:} a doctor or fracture clinic nurse checks the fit, comfort, and circulation of the limb
    \item \textbf{Neurovascular checks:} testing colour, warmth, feeling, and movement in the exposed hand or foot
    \item \textbf{Splitting or removal of the cast:} if it is too tight, it may be bivalved or removed to relieve pressure
    \item \textbf{X-ray:} used to confirm that the fracture is still aligned and healing as expected
    \item \textbf{Skin inspection:} if irritation is suspected, a window may be cut in the cast to examine the skin underneath
\end{itemize}

\section*{Management}
\begin{itemize}
    \item \textbf{Keep it dry:} cover the cast with a waterproof sleeve or plastic bag when showering or bathing; never submerge it
    \item \textbf{Elevate:} raise the injured limb on pillows above heart level for the first few days to reduce swelling
    \item \textbf{Apply cold:} an ice pack placed over the cast (never directly on the skin, and only when the cast is dry) can ease discomfort
    \item \textbf{Move freely:} regularly wiggle the fingers or toes and move the uninjured joints to keep them mobile
    \item \textbf{Do not poke or scratch:} never insert objects such as knitting needles or coat hangers down the cast
    \item \textbf{Protect the cast:} avoid weight-bearing or activities that the doctor has advised against
    \item \textbf{Follow the plan:} attend scheduled reviews and cast removal, and follow instructions about when weight-bearing is allowed
    \item \textbf{Muscle care:} once the cast is removed, gentle exercises and physiotherapy help restore movement and strength
\end{itemize}

\section*{Complications}
\begin{itemize}
    \item \textbf{Compartment syndrome:} dangerously high pressure within the limb muscles, a medical emergency that can cause permanent muscle and nerve damage
    \item Pressure sores and skin ulcers where the cast rubs or presses
    \item Skin infection, maceration, or fungal infection if moisture becomes trapped under the cast
    \item Nerve compression causing numbness or weakness, usually reversible if the cast is loosened early
    \item A wet or damaged cast that fails to hold the bone, leading to delayed or failed healing
    \item Joint stiffness and muscle wasting from immobility, usually improving with exercise after removal
    \item Itching, heat rash, or dermatitis from the cast material
\end{itemize}

\section*{Prognosis and Outcomes}
Most people heal well with a cast and return to normal activity once it is removed. Simple fractures typically take several weeks to heal, and the treating team decides when the cast can come off based on X-rays and the expected healing time. After removal, the limb may be stiff, weak, and slightly swollen for a time, and physiotherapy speeds full recovery. Complications are uncommon when the cast is cared for properly and warning signs are reported early; compartment syndrome is the main reason prompt reporting of severe pain is important.

\section*{Prevention and Self-Care}
\begin{itemize}
    \item Keep the cast completely dry at all times; use a waterproof cover or bag in the shower
    \item Elevate the limb and move the fingers or toes early and often to control swelling
    \item Never put anything inside the cast to scratch; use a hairdryer on a cool setting for itching
    \item Do not trim, reshape, or remove the cast yourself -- leave this to the clinic
    \item Avoid heat: do not use heaters or hairdryers to dry a wet cast, as heat can damage the skin
    \item Report severe pain, numbness, colour change, or a bad smell without delay
    \item Follow all weight-bearing and activity instructions given by the treating team
    \item Attend follow-up appointments for review and cast removal
\end{itemize}

\section*{Sources and Further Reading}
The following organisations provide reliable information on fracture care, casting, and cast complications.

\begin{itemize}
    \item National Health Service. \textit{NHS conditions guide on broken bones and cast care.} https://www.nhs.uk/conditions/
    \item Mayo Clinic. \textit{Information on casts, splints, and immobilisation care.} https://www.mayoclinic.org/
    \item MedlinePlus. \textit{Consumer health information on casts and cast removal.} https://medlineplus.gov/
    \item MSD Manual. \textit{Overview of fracture management and casting.} https://www.msdmanuals.com/
    \item Centers for Disease Control and Prevention. \textit{Injury prevention and fracture information.} https://www.cdc.gov/
    \item World Health Organization. \textit{Global guidance on musculoskeletal injury care.} https://www.who.int/
\end{itemize}
